% Shared property/database workflow, inspected 12 September 2026.
% New displays are unnumbered to preserve the mathematical equation sequence.
\subsection{Independent precision and replacement rules}
The database stores one set of physical properties per theory. A Lagrangian
realization supplies the input for later calculation; different realizations
attached to the same theory share its indices and Coulomb spectrum.
Initial imports run the anomaly/conformality checks and calculate basic
properties only. Both indices, the spectrum and both precision columns start
as SQL \code{NULL}; their keys are initially absent from \code{properties_json}.

The separate calculator returns full-index precision $N\in\mathbb Z_{\geq0}$
and Coulomb precision $D\in\mathbb Q_{\geq0}$. These are inclusive cutoffs:
\[
 \mathcal I_{\leq N}=\sum_{a=0}^{N}t^a P_a(y,u),\qquad
 \mathcal I_{\mathrm C,\leq D}=\sum_{0\leq\Delta\leq D}c_\Delta x^\Delta.
\]
A zero boundary coefficient does not reduce the trusted cutoff. For example,
$1+x^2+x^4$ computed through $D=5$ must retain precision five, even though its
largest nonzero exponent is four. The complete gauge-invariant generator
spectrum is independent of either cutoff; it is not truncated with the index.

For each of the two indices separately, let $V$ be its stored value and
$p_{\rm old},p_{\rm new}$ its stored and submitted cutoffs. The write rule is
\[
 \text{replace or fill}\quad\Longleftrightarrow\quad
 V\text{ is missing}\quad\text{or}\quad
 \bigl(p_{\rm old}\text{ is known and }p_{\rm new}>p_{\rm old}\bigr).
\]
\begin{tabularx}{\textwidth}{@{}p{50mm}Y@{}}
\toprule Stored / submitted state & Database action\\\midrule
Stored index missing & Fill it and record the supplied cutoff.\\
Known, strictly higher cutoff & Replace that index and its cutoff.\\
Equal or lower cutoff & Keep the existing string and precision unchanged.\\
Legacy index, unknown cutoff & Keep it; report \code{unknown_precision}.\\
Omitted or \code{None} component & Leave the corresponding stored data alone.\\
Supplied index without a cutoff & Reject the payload before opening the write transaction.\\\bottomrule
\end{tabularx}
For example, an update from $(N,D)=(18,90)$ to $(24,60)$ upgrades the full index
only. The Coulomb index remains at dimension 90. Equal-order submissions do
not replace even a different string; this API is not a correction/force-update
mechanism. The caller supplies the trusted calculation cutoff; the writer does
not infer or prove precision from the expression's last monomial.

A missing spectrum is filled. An identical complete spectrum is retained, with
multiplicities preserved; a conflicting spectrum raises an error and rolls back
that update call. This policy differs from comparison of truncated indices.

{\small\textbf{Implementation source:} \code{common/n2_theory_properties.py},
\code{common/n2_theory_db.py}; the inequalities above express the implemented
cutoff contract, not an additional physical conjecture.}

\clearpage
\subsection{Schema 7, legacy data and atomic writes}
Schema 7 keeps the same ten tables and adds two nullable columns to
\code{theory_properties}. Central charges and all existing index data are
preserved by the migration.
\begin{tabularx}{\textwidth}{@{}Y>{\raggedright\arraybackslash}p{57mm}@{}}
\toprule Column / result key & Stored representation\\\midrule
\code{superconformal_index_order} & \texttt{BIGINT UNSIGNED}; inclusive integer $N$.\\
\code{coulomb_branch_index_max_dimension_json} & Exact numerator/denominator JSON for $D$.\\
\code{coulomb_branch_index_max_dimension} & Corresponding key in calculator output and combined JSON.\\\bottomrule
\end{tabularx}
The three computed values still use \code{superconformal_index_json},
\code{coulomb_branch_index_json} and \code{coulomb_branch_spectrum_json}.
Successful updates mirror changed values and cutoffs into \code{properties_json}
while preserving basic properties and other computed components.

Migration $6\to7$ adds the columns without guessing legacy cutoffs, recomputing
indices, rehashing realizations or merging theories. Initialization applies
migrations when the database is next opened through the database API. MySQL
schema DDL precedes the explicit data transaction and is not rolled back with it.

When an old result's original requested cutoff is independently known, record
it before normal upgrades:
\begin{lstlisting}[language=Python]
from common.n2_theory_db import record_lagrangian_index_cutoffs

# Use these numbers only if they are the verified original cutoffs.
record_lagrangian_index_cutoffs(
    connection, realization_id, order=18, max_dimension=90,
)
\end{lstlisting}
This function fills unknown precision only. It rejects changing an already
known cutoff or annotating an index that is absent. It changes neither index
expression. Merely seeing a last term of degree 18 does not establish that the
original calculation was requested through 18.

\code{update_lagrangian_indices(connection, realization_id, indices)} accepts
\code{calculate_n2_theory_indices}'s output or a subset of computed components
with their matching cutoffs. It opens a short transaction, locks the shared
row with \texttt{SELECT ... FOR UPDATE}, rechecks the current cutoffs, and merges
accepted changes into the columns and JSON together. Failure rolls back the
call; a no-op does not change the theory's \code{updated_at} timestamp.

Calculations run before this lock is acquired. Thus a slower low-order worker
cannot overwrite a higher-order result committed while it was calculating.
Use an initialized connection with no caller-owned transaction for update and
legacy-cutoff APIs. These APIs manage their own commit/rollback boundary.

\clearpage
\subsection{Separate worker, partial retries and reproducible use}
The program \code{common/n2_theory_db_indices.py} streams stored Lagrangian
inputs in bounded pages, using the smallest realization ID for each theory.
It visits each shared theory once rather than recalculating the same shared
properties for every attached realization.
\begin{enumerate}
\item Default mode selects any missing full index, Coulomb index or spectrum.
SQL \code{NULL} and JSON \code{null} are both treated as missing.
\item \code{--upgrade} additionally selects indices with known lower cutoffs.
Already sufficient components are not calculated. Complete legacy records with
unknown precision are reported without recalculating those indices.
\item Each requested component is calculated outside a database transaction,
then submitted to the locked update API. The stored state is checked again at
write time, even if it differed when the job was selected.
\item Each successful component commits independently. A failure is reported
for that component; subsequent components and theories continue. Re-running
retries remaining missing or lower-order components and keeps earlier successes.
\end{enumerate}
Multiple workers can duplicate a calculation; there is no persistent job-claim
queue. The row-lock and cutoff policy protects data against downgrades, not
against duplicated CPU work. A conflicting spectrum rolls back its own update
call, not components already saved by earlier worker calls.

Run from the repository root with Sage's Python. Database credentials follow
the existing \code{N2_DB_HOST}, \code{N2_DB_USER}, \code{N2_DB_PASSWORD} and
\code{N2_DB_UNIX_SOCKET} environment conventions; port is a CLI option.
\begin{lstlisting}[language=bash]
# Stage 1: validate and store basic properties.
sage -python -m common.n2_theory_db \
  anomalies/example_e6.json DATABASE --user USER
# Stage 2: fill missing components, defaults N=18 and D=90.
sage -python -m common.n2_theory_db_indices DATABASE --user USER
# Later: upgrade only components with lower known precision.
sage -python -m common.n2_theory_db_indices DATABASE --user USER \
  --upgrade --index-order 24 --coulomb-max-dimension 100 --limit 10
\end{lstlisting}
The worker also accepts \code{--host}, \code{--port}, \code{--unix-socket},
\code{--connect-timeout}, \code{--cache-directory}, \code{--lie-executable},
\code{--form-executable}, \code{--timeout} and \code{--processes}.
\code{--limit} bounds visited theory jobs, including reported legacy skips.
The API \code{fill_lagrangian_indices} yields the same per-theory outcomes.
Each outcome reports \code{updated_fields}, \code{skipped_fields} and
\code{errors}; a final CLI summary counts processed and failed jobs. Exit codes
are 0 for no component failures, 1 for component failures, and 2 for an outer
configuration or database-level error.

\textbf{Validation recorded during implementation, 12 September 2026.}
All 238 backend tests passed in 21.920 seconds, including an isolated MySQL
8.0.46 server, schema migration, concurrent upgrades, rollback, partial retries,
JSON-null selection and real Sage/FORM/LiE calculations. This PDF update checks
source contracts and document rendering; it does not rerun that suite or alter
the user's database. Older cache-performance measurements remain dated evidence.
